This manuscript was developed through an extended human--AI research
workflow.  We disclose that workflow because the distinction between setting
the scientific problem, generating candidate arguments, and validating a
theorem is material here.

\paragraph{Baojian Zhou's contribution.}
Baojian Zhou conceived and directed the project, fixed the PageRank and local
work questions, and made the scientific scope decisions.  In particular, he
identified that the canonical theorem should use the point source
$\vs=\eunit{v}$ rather than silently claim the stronger general sparse-source
bound; proposed the discover-once/solve-once AESP--CD viewpoint; recognized
the importance of the 2023 COLT active-set bottleneck and the subsequent
Wei--Yang result; identified the grounded-flow dual as a potentially
important organizing viewpoint; selected which claims, counterexamples, and
proof routes should be retained or rejected; and requested the literature,
notation, and complexity audits that shaped the present theorem.  He is the human author,
made the final research and exposition decisions, and takes responsibility
for the manuscript.

\paragraph{GPT-5.6-sol assistance.}
OpenAI's GPT-5.6-sol model, used through the Codex research and coding tools,
served as an interactive research assistant during the 2026 proof-development
and drafting period.  With access to the local repository and author-supplied
paper archive, it helped search and reconcile notation, trace statements to
primary sources, derive and stress-test candidate inequalities, construct and
check small counterexamples, organize claim ledgers, and draft and revise
LaTeX.  In the central proof, the tool assisted in turning the author's
repeated-face bottleneck into the threshold-batch formulation, developing and
auditing the block-Cholesky/Chebyshev decay route, checking the fully charged
SDD wrapper, and formalizing the approximate-inner-face two-stage transfer.
It also ran repository tests and manuscript builds and reported the resulting
diffs for human review.

\paragraph{Validation and responsibility.}
AI-generated text or algebra was never treated as mathematical evidence.
Candidate claims were kept only when they had an explicit proof, a
source-level citation, or a clearly labeled computational audit; failed and
conditional routes remain separated in the repository notes.  Literature
claims were checked against the cited primary papers rather than against model
recollection.  The AI system is not an author, cannot approve the final
claims, and bears no responsibility for errors.  The human author is
responsible for independently checking every theorem before submission and
for the final content, citations, and conclusions.
